\documentclass[final,3p,times,11pt,authoryear]{elsarticle}

% Material suplementar: mesmo preambulo do artigo principal. Gerado por make_latex.py a partir de supplementary_material.md.

\usepackage[utf8]{inputenc}
\usepackage[T1]{fontenc}
\usepackage{amsmath,amssymb,amsfonts}
\usepackage{graphicx}
\usepackage{textcomp}
\usepackage{booktabs}
\usepackage{multirow}
\usepackage{url}
\usepackage{float}
\usepackage{tabularx}
\usepackage{setspace}
\onehalfspacing
\usepackage{enumitem}
\setlist{nosep,topsep=2pt,itemsep=1pt,parsep=0pt,leftmargin=*}
\usepackage{xcolor}
\usepackage{hyperref}
\hypersetup{colorlinks=true,linkcolor=blue!50!black,citecolor=blue!50!black,urlcolor=blue!50!black}
\emergencystretch=3em
\newcommand{\todo}[1]{\textcolor{red}{\textbf{[TODO #1]}}}
\renewcommand{\thetable}{S\arabic{table}}
\renewcommand{\thefigure}{S\arabic{figure}}

\journal{European Journal of Operational Research}

\begin{document}
\begin{frontmatter}
\title{Supplementary material to: %%TITLE%%}
\author{Bonfim Amaro Junior, Marcio Costa Santos, Guilherme Nepomuceno de Carvalho, Plácido Rogério Pinheiro, Carlos Vinicio Dantas Araújo}
\begin{abstract}
This document contains supplementary material to the article: details of the attempts to repair adaptive operator selection, average ranks by instance family, and the mechanism of the effect of a tight cost budget. References are those of the main article.
\end{abstract}
\end{frontmatter}

%%BODY%%

\begingroup
\scriptsize
\setlength{\bibsep}{0pt}
\bibliographystyle{elsarticle-harv}
\bibliography{referencias}
\endgroup

\end{document}
